\chapter{Índice general y de recursos complementarios}
La tabla de contenidos o índice es una parte esencial del informe ya que sirven de guía para la navegación a través del contenido. Además del índice, se recomienda la inclusión de un glosario de términos, lista de abreviaturas, símbolos, figuras y tablas. A continuación, se presenta una descripción de cada uno de estos componentes.

\section{Índice}
Corresponde a una lista de las principales divisiones y secciones del texto en el mismo orden en que aparecen, indicando las respectivas páginas iniciales. El índice debe ser ubicado inmediatamente después de la portada y agradecimientos en el caso de un TFT o TFG. 

\section{Glosario de Términos}
Corresponde a una lista en orden alfabético de palabras especiales de significado poco conocido o palabras en otro idioma, acompañada de sus respectivas definiciones. Cuando se incluye, el glosario debe ser situado justo antes de la Introducción.

\section{Índices de Recursos Complementarios}
Además de la tabla de contenido y el glosario de términos, en informes extensos es recomendable incluir índices adicionales para proporcionar una referencia rápida y detallada a elementos específicos y técnicos utilizados a lo largo del documento. Estas listas adicionales suelen incluir:
\begin{itemize}
    \item Lista de Símbolos: Define los símbolos poco conocidos utilizados en el orden en que aparecen en el texto, acompañado de sus respectivos significados.
    \item Lista de Figuras/Tablas/Algoritmos: Se presentan en el orden en que aparecen en el texto, indicando para cada una de ellas su índice, el cual está compuesto por el número del capítulo y el número de aparición en el capítulo. Además, se debe indicar nombre y página en donde se encuentra.
\end{itemize}